\section{Critical-aspect angular signs with diverging transverse row mass}\label{sec:critical-aspect-angular-nodes}

Keep the notation of \href{https://github.com/ProgramComputer/open-math-research/blob/4fb624ba937a8eecc1361a6066185cd81ab7b172/topics/balanced-arrays/quadratic/very-tall-quadratic-interior-saddle.md}{the quadratic interior saddle}, with $1\le r<d$, $d\to\infty$, and assume $n/d^2$ stays in a fixed compact subset of $(0,\infty)$. Suppose $\mu\to\infty$ and $d-\mu\to\infty$. Then $da\to\infty$: indeed $a=[\mu(d-\mu)-V]/[d(d-1)]\ge\mu(d-\mu)/d^2$, using $V\le\mu(d-\mu)/d$. In particular $na/d\to\infty$. All the angular estimates below only require $da\to\infty$, $\mu\ge1/10$, and the stated compact aspect bounds.

Write $D=\lfloor nr/d\rfloor$, $t_*=\pi/d$ and

\begin{equation*}
\Psi=n\arg\phi(t_*)-\tfrac{d-1}{2}\arg a(t_*).
\end{equation*}

Take the argument of $\phi(t_*)$ continuously in the real tilt (equivalently, the argument in $[0,r t_*]$), and the argument of $a(t_*)$ near zero. At any nodes $\Psi=k\pi$, $1\le k\le D-1$, one has $(-1)^kP_n^{(r,d)}(-v^d)>0$ eventually along every sequence satisfying the assumptions. No lower-quota estimate is used. Fixed corank with bounded $d-\mu$ is outside the asserted range unless $da\to\infty$ can be checked separately.

\subsection{A quantitative single-chart estimate}

For $h=r-\mu\ge h_0$ with fixed small $h_0>0$, use the near charts $|t|\le c_0/\sqrt{1+V}$ from the quadratic proof and let $b_t=(1+V)t^2$. For $0<h\le h_0$ use all charts and let $b_t=h\min(t^2,1)$, using the diagonal and determinant bounds in \href{https://github.com/ProgramComputer/open-math-research/blob/4fb624ba937a8eecc1361a6066185cd81ab7b172/topics/balanced-arrays/quadratic/near-full-all-angular-signs.md}{the near-full angular analysis}. In both cases $\phi(t)$ is bounded away from zero in the respective range. With normalized Haar charts and $I_0=(2\pi n a)^{-(d-1)/2}/\sqrt d$, write

\begin{equation*}
G_t=I_0\, \phi(t)^n\left(\frac{a}{a(t)}\right)^{(d-1)/2},\qquad
e_d=\sqrt{\frac d{na}}+\frac d{na}+\frac{d^{5/2}}{n^{3/2}}.
\end{equation*}

The integral of the normalized row response $R^n$ over each such chart is
\[G_t+O(I_0 |\phi(t)|^n e^{Cdb_t}e_d).\] The modulus of the full chart integral is at most $C I_0|\phi(t)|^n e^{Cdb_t}$. Constants can depend on the compact aspect bounds and $h_0$.

To prove the estimate, use $Q_j=\sum y_i^j$ and $L=\sum_{i\in I}y_i$. The log centered chart response has exactly quadratic coefficient $-a(t)Q_2/2$, exactly cubic coefficient $-i\gamma(t)Q_3/6$, with $|\gamma(t)|\le Ca$, and a fourth log coefficient bounded by $Ca Q_4$, by \hyperref[sec:fourth-transverse-cumulant]{the transverse cumulant calculation}. The latter proof uses only bounded $\phi(t)^{-1}$, $a\ge V/d$ and nonnegative radial weights, so it applies to every near-full phase as well. The error beyond these terms is bounded by $C\mathbb E|L|^5$ for all $\|y\|\le\eta$ with fixed small $\eta$. In detail, $\delta=R(t\mathbf1+y)/\phi(t)-1$ has $|\delta|\le C\mathbb E|L|^2$ by its zero linear term, and the Taylor expansion of $\delta$ has raw remainder at most $C\mathbb E|L|^5$. In $\log(1+\delta)$ only the square of its quadratic term contributes through order four. The rest of the logarithm is bounded by $C[(\mathbb E|L|^2)(\mathbb E|L|^3)+(\mathbb E|L|^2)^3]$, which is at most $C\mathbb E|L|^5$ by monotonicity of moment norms and $\mathbb E|L|^2=O(\eta^2)$. This estimate uses the bound on $|\delta-\delta_2|$ directly rather than squaring the order-four Taylor polynomial.

Under the real Gaussian of precision $A$ comparable to $na$ on $\mathbf1^\perp$, the fixed-support $L$ has variance $s(d-s)/(d A)=(d-1)a_s/A$. Minkowski with respect to the subset law gives

\begin{equation*}
\|\mathbb E|L|^5\|_{L^2(\mathrm{Gauss})}
\le C(d/A)^{5/2}\mathbb E a_S^{5/2}
\le C(d/n)^{5/2}.
\end{equation*}

For the last bound use $O(1/d)$ Lipschitz continuity of $a_s$, $a_\mu=a+V/[d(d-1)]$, $a\ge \max(V/d,1/(11d))$ and $\mathbb E|S-\mu|^{5/2}\le C(1+V)^{5/4}$. That moment bound follows from the log-concave survival estimate of the quadratic proof, which holds without a positive lower bound on $V$. The cubic and canceled quartic terms have Gaussian $L^2$ norms $O(\sqrt d\,a A^{-3/2})$ and $O(d a A^{-2})$ by Wick and the Gaussian eighth moment. Multiplying by $n$ gives $O(e_d)$.

For exponential domination in Cauchy--Schwarz, retain the original crude real quartic bound from the quadratic and near-full proofs, namely $CnaQ_4+Cna^2Q_2^2$ with the real cubic treated as in those proofs. In the radial $\chi^2$ central region it costs only $\exp(O_L(d^2/n))=O(1)$ for fixed large $L$ and aspect bounded away from zero; the tail is dominated by a weakened Gaussian just as before. The coordinate bounds cost at most $e^{C d b_t}$ (use the tilted coefficient derivatives from the near-full proof, or Young's inequality in the quadratic proof). Thus the difference between the $n$th powers of the exact chart response and of the complex Gaussian is controlled by the stated bound, by $|e^x-e^y|\le |x-y|e^{\max(\mathrm{Re}\,x,\mathrm{Re}\,y)}$ and Cauchy--Schwarz. The Gaussian tail outside each fixed-radius chart is negligible, since $na\gg d$. The exact Gaussian gives $G_t$, using torus covolume $\sqrt d(2\pi)^{d-1}$. Here $e_d=o(1)$ uniformly along any admissible sequence.

The complement of the charts is negligible on the principal scale by the quadratic proof's integrated relative-angle bound. Its relative bound is $\exp(O(d)-cna+O(nV/d^2))=o(1)$ on the unsaturated chord region, since $V/d^2\le a/d$ and $na\gg d$. The saturated chord region is bounded by $\exp(-cn)$ before normalizing, which absorbs both the principal radial loss and $O(d\log n)$. For $h\ge h_0$ the non-near diagonal charts also have the constant row modulus gap shown there and are negligible.

\subsection{Phase sum away from saturation}

Here $V\ge c_{h_0}$. If $V\to\infty$, the radial gap in the quadratic proof and the uniform chart bound make all other phases negligible relative to one principal saddle: since $nV\gg d(1+V)$, the near phases with $|t|=j\pi/d$, odd $j\ge3$, decay as $\exp[-c nV(j^2-1)/d^2]$. The two principal terms at a node then have the same real sign $(-1)^k$.

On any subsequence with $V$ bounded above, the same bound makes the sum over $j>M$ arbitrarily small on the principal scale by taking $M$ sufficiently large, depending only on $h_0$ and the aspect bounds. For each fixed signed odd $j$, use
$\arg\phi(t)=\mu t+O(|t|^3)$, since the centered third absolute moment is uniformly bounded. Also $\arg a(t)=\xi t+O(t^2)$ for a real $\xi$, with uniform constant, by the weighted second-moment estimates in the quadratic proof. Thus the phase of $G_{j t_*}$ differs from $j\Psi$ by $O_M(n/d^3+1/d)=o(1)$. Their real parts all have sign $(-1)^k$ eventually, and include the principal magnitude. The sum of their local approximation errors is smaller by $o(1)$. This proves the sign in the bounded case and then by subsequences for all $h\ge h_0$.

\subsection{Phase sum near saturation}

Here $V\asymp h$ and $\log|\phi(t)/\phi(t_*)|\le-ch(t^2-t_*^2)$ for $t_*\le |t|\le\pi$, by \href{https://github.com/ProgramComputer/open-math-research/blob/4fb624ba937a8eecc1361a6066185cd81ab7b172/topics/balanced-arrays/quadratic/near-full-all-angular-signs.md}{the near-full estimates}. The same estimate on node positions there gives $T-k=(nh/d)(1+O(d^{-2}+d/n))$ with $T=nr/d$, so $nh\gtrsim d$ at every internal node. Choose $M$ large and retain signed odd phases $|t|\le M/\sqrt{nh}$, including the principal pair since $nh/d^2$ is bounded above. Write $\ell=t/t_*$, an odd integer. The argument of $H(ve^{it})/H(v)$ may be lifted as $\mu t+O(h|t|^3)$. Also $\arg a(t)=\xi t+O(ht^2)$ for a real $\xi$, because the first and second derivatives of $a(t)/a$ are $O(h)$ uniformly. Subtracting $\ell$ times the principal phase of $G$ cancels the two linear terms. Since $|\ell|t_*^j\le |t|^j$ for $j=2,3$, the phase difference is

\begin{equation*}
O(nh|t|^3+dh t^2)=O_M((nh)^{-1/2}+d/n)=o(1).
\end{equation*}

The retained real parts all have sign $(-1)^k$. Their total amplitude on the principal scale is at least $c(1+d/\sqrt{nh})$ by counting the odd phases in $|t|\le1/\sqrt{nh}$ when that count is large, and otherwise using the principal pair. Their local errors are a relative $o(1)$ because $db_t\le M^2d/n=o(1)$. As in the earlier near-full proof, the discarded total is at most $C(1+d/\sqrt{nh})e^{-c'M^2}$ on the principal scale, by the radial phase gap, absorption of $Cdh t^2$ since $n\gg d$, and the resulting Gaussian lattice tail. Taking $M$ sufficiently large gives the sign at every near-full internal node.
